\section{Necessary restrictions, verification, and comparison}\label{sec:scope}
\subsection{What the assumptions do and do not assert}
The strict normal-cone condition is sufficient for the lifting theorem; we do not claim it is necessary for every form of local-search hardness. It is essential to this confinement mechanism. If $B_0\notin\mathcal V_H$, separation gives a leading tangent direction $X$ with $\ip{B_0}X<0$. At any fixed vertex of $\mathcal F_H$, extend it to a full tangent direction by choosing its bottom entry sufficiently large: all finitely many positive-last-coordinate active inequalities can then be satisfied, while the leading active inequalities already hold. Decomposing the pointed tangent cone into its extreme rays, at least one ray has negative leading cost. For large $N$, that fixed ray improves the scaled objective and leaves the face.

Boundary exposure can also fail. At $T_3$, the direction
\[
 R_{12}=R_{21}=1/2,\qquad R_{13}=R_{31}=-1/2,
\]
with all other entries zero is an extreme tangent ray: among the six interval inequalities only the interval $[1,2]$ has positive value. With $B=I_3+(e_1e_3\trans+e_3e_1\trans)/4$, the leading block $I_2$ lies on the boundary of $\mathcal V_{T_2}$ and $\ip{S_NBS_N}R=-N/4<0$. This example diagnoses the exposure condition; it is not claimed to be a nonmember lower-bound example.

Multiple starts sharing the leading block are covered in full. Arbitrary menus of distinct leading perfect blocks are not covered. If $H_1\ne H_2$ are distinct perfect vertices, the interiors of their minimizing normal cones are disjoint: a functional in either interior uniquely minimizes at its vertex. Thus one leading seed cannot satisfy both strict exposure conditions. This is an obstruction to this proof strategy, not a theorem that all diverse-start methods are efficient or that no other hardness construction exists.

The extreme-ray restriction is also real. Adding a sufficiently small positive multiple of an escaping feasible tangent direction to a strictly improving face ray still gives a feasible improving direction, but leaves the face. Consequently this proof does not cover arbitrary nonextreme descent directions, nonlocal jumps, or input-dependent initializations. Arbitrary global certificate generation is intentionally outside the lower bound; $W_N$ is a short explicit global separator.

Positive definiteness is needed for the independent geometric theorem. For example $H=0$ gives $\mathcal E_H=\{b\ge0,c\ge1\}$, which has only one vertex rather than polynomial graph growth of positive degree. All algorithmic claims with fixed $H,n$ allow constants depending on both; the template bounds can be very large and do not imply uniformly practical computation in variable dimension.

\subsection{Relationship to prior work}
The simplex model, perfect matrices, finite normal cones, and qualitative Ryshkov geometry are antecedents \cite{DSV,DS,OS}. The bounded-edge displacement estimate and the canonical five-cycle all-rule construction were proved in the preceding project manuscript, and are reproved here rather than relabeled as new. The strengthening is the finite-template normalization and its sharper ray bounds, the seed-based linear-scale lifting theorem, its stronger certificate exclusions, higher-order principal-minimal applications, and the independent graph-growth theorem.

The PSD-principal-block lemma is explicitly attributed to \cite{SM}; it is not a new matrix decomposition. Odd-cycle obstructions and the distinction between doubly nonnegative and completely positive matrices are classical. The new odd-cycle conclusion couples these seeds to all-choice canonical hardness and retains complete positivity of every proper principal block, even on a specified neighborhood. Motzkin--Straus supplies the classical graph witness; a self-contained proof was given above. The unrestricted lattice geometry is related to \cite{Erdahl,IE}; our conclusions concern a different, orthant-truncated rational epigraph and its full adjacency graph. The recent algorithmic work \cite{Minimum} studies the copositive minimum rather than a lower bound on adjacent traversal.

These comparisons isolate the proposed new assertions. They are not an exhaustive first-priority certificate. A lack of an identical result in the sources inspected is not a proof of historical novelty. The earlier suitable-pivot finite-termination conjecture is not contradicted: an exponentially long path may still terminate. We do not assert eventual termination of every permitted local rule.

\subsection{Reproducibility and status}
The supplement contains exact rational seed certificates, finite normal-system checks, shear-template checks, proper-principal CP factors, and the two-dimensional graph model. A separately written verifier imports no generator code and rejects deliberately corrupted certificates. The original canonical certificate suite is retained and rerun without modifying its formulas. The verification report records the actual case counts and file hashes; finite tests supplement, rather than replace, the proofs.

No production high-dimensional polyhedral library, external human proof review, or proof-assistant formalization is claimed. The user-reported randomized ten-paper benchmark is not recalibrated in this mathematical supplement: the sample and its complete comparison rubric were not available in the retrieved research files. The revised contribution is suitable for a fresh skeptical comparison, but the paper does not assert a top-three rank or an editorial probability.
